\chapter{What controls motion on a shared road?}
\label{ch:g8}
The locked study measured the behavior of the complete synthesis system. It left two questions open. Does the diffusion network learn a class-dependent motion law when road and noise are fixed? Can the inherited image score improve that law by acting during denoising? The class intervention below uses development participants. The feedback experiments then move from development cases to two independently simulated cohorts, with rules fixed before inspecting each new cohort. The distinction matters because a preference observed while designing a selector can fail on fresh traffic.

\section{A class switch on the same road}
We took the first planned path for each of the eight bus and eight car development requests. Each full path passed the directed graph audit under both class permissions. We then fed each path and start to each of the three frozen diffusion models. Each of eight initial Gaussian noise draws was used twice, once with a bus label and once with a private-car label. The deterministic sampler, path geometry, checkpoint and noise were identical within each pair. No GAF selection was performed. The 768 decoded traces all passed the graph audit and none exhausted its path.

\fig{g8-switch.pdf}{1}{Change in generated motion when only the input class changes. Gray lines represent 16 fixed development routes, averaging eight noise draws and three fixed models. The darker line joins the means.}{switch}
Mean travel was 984.3 metres with the bus label and 1,330.4 metres with the car label. Bus-minus-car differences for the three weights were $-319$, $-375$ and $-344$ metres. The paired mean difference was $-346$ metres, with a conditional 95\% interval from $-469$ to $-213$ metres. The bus label raised the proportion of intervals below 0.1 m/s from 5.25\% to 23.83\%. The difference was 18.59 percentage points, interval $[11.89,25.34]$. Fourteen of the sixteen routes had a positive mean stop-fraction difference. Intervals use 2,000 whole-route resamples stratified by original request class, conditional on fixed draws and weights. Seed 42 motivated the subsequent checks with seeds 43 and 44 on the same cohort.

Observed development buses traveled 948 metres on average, versus 1,340 metres for cars, with near-stationary fractions of 29.7\% and 3.7\%, respectively. The signs agree with the class switch. The observed groups drove different routes and provide no ground truth for what a real bus would do in place of a particular car. The experiment identifies the response of this model to a label intervention. It also identifies class information available to a different representation: the independent trajectory auditor's mean bus probability rose from 0.230 to 0.444 when the label changed from car to bus.

\section{Does an image detect motion when geometry is fixed?}
For each of the 33 development requests, we kept the first seed-42 path, its starting position and its total arc travel. We resampled cumulative progress at $u=t+a\sin(2\pi t)$ for normalized time and $a\in\{-0.12,0,0.12\}$. These monotone time maps preserve the endpoints. All 33 requests produced a different image under at least one perturbation; every variant was encodable. Selecting the highest-scoring of the three variants increased the intended-class image probability by 0.257 on average relative to their within-request mean. The independent TCN probability changed by $-0.004$. Mean discrepancy from the observed development class speed distribution increased by 0.289 m/s. These are per-request descriptive comparisons, with one parent motion per request. The scorer detects timing even on fixed geometry, but its preferred timing did not improve the independent measurements in this probe.


\section{Why recognition can succeed while generation gains little}
The road and requested class already condition each generated sample. A further classifier score rewards a different property. Consider an ideal calibrated classifier on a complete trace $x$ and an ideal conditional generator $p(x\mid c)$. Weighting samples by $p(c\mid x)$ favors traces that distinguish class $c$ from the other classes. By Bayes' rule, this factor depends on the ratio $p(x\mid c)/p(x)$, whereas motion fidelity is assessed against $p(x\mid c)$ itself. For example, many ordinary car and taxi movements may look similar. A taxi scorer can assign a higher score to an unusual, easily recognized motion than to the common motions that a synthetic taxi population ought to contain. Even a perfect recognizer would therefore need an independent empirical test of its usefulness for generation. Our finite GAF scorer uses an image of the trace and is imperfect. This argument explains the distinction between the objectives without assuming an exact density for our resampling algorithm.

Within the initial candidate experiment, GAF detected timing changes on a fixed path. Introducing a single scoring decision during denoising also changed the dispersion of the returned candidates. On higher-turn development roads, one intervention at update 76 appeared to bring speeds closer to observations. That speed advantage failed to replicate on eight fresh SUMO traffic episodes. Details of the timing controls, multiple outcomes and the first independent replication are retained in Appendix~\ref{app:fixedtime}. These findings raise a more direct question: does scoring help when the current candidate motions have become distinguishable, irrespective of a hand-picked update number?

\section{Waiting for distinguishable motion}
For each candidate $k$, let $s_k^{(j)}$ denote the 512-point cumulative progress profile decoded from its current clean prediction, capped at its planned path length. Every five reverse updates we measure the typical change in a candidate and the typical separation between candidates:
\begin{align}
 D_j &= \operatorname{median}_{k}\;
       \operatorname{RMS}\big(s_k^{(j)}-s_k^{(j-5)}\big),\\
 B_j &= \operatorname{median}_{k<l}\;
       \operatorname{RMS}\big(s_k^{(j)}-s_l^{(j)}\big).
\end{align}
Both quantities are in metres. We score once when a second consecutive check first satisfies $D_j<B_j$ with $B_j>0$. Thus a single lucky fluctuation cannot set the intervention time. The five-update cadence determines measurement resolution and was fixed before this cohort. If the criterion never fires before the final transition, the policy uses terminal selection. This heuristic examines motion geometry and uses no observed future, image scores or validation loss to trigger. Relative stability need not imply that the image classifier's eventual ranking is correct.

At the chosen update, the pipeline converts the eight provisional decoded motions into the inherited Hilbert images and uses the intended-class probabilities for systematic resampling with an effective-sample-size floor of six. A stochastic transition gives the children distinct continuations. The terminal comparator uses the same transition, paths, noise, score budget and weight rule but scores the finished trajectories. Each logical policy uses 100 denoising updates per particle and eight images per decision. Computing the progress diagnostic adds overhead. For experimental reproducibility the trigger-discovery pass was replayed; equal denoiser budgets here describe the policies, not a wall-clock speedup of the implementation.

To identify the contribution of matching the time to a particular pool, we also permute trigger times among requests of the same class, separately by model and noise repetition. The full histogram of intervention times is unchanged. The paired effect of the permuted policy is calculated against its own transition-matched terminal comparator. Subtracting these two paired effects tests the value of assigning a particular time to a particular state. The permutation is an offline experimental control.

\section{Motion quality on a fresh cohort}
We generated eight further SUMO demand and traffic seeds (260201--260208) on the same MoST map, following the independent simulation protocol in Appendix~\ref{app:fixedtime}. Of 800 target vehicles, 416 met the strict reconstruction criteria; the predeclared quotas retained 240 requests. The primary endpoint was speed-distribution discrepancy across all five classes. Three existing diffusion models and three new noise draws yielded 2,160 returned traces per policy, with whole-episode paired resampling for uncertainty.

\fig{g11-trigger.pdf}{.92}{The proposed rule generated a distribution of stopping times across the 2,160 pools. The median was update 30, with observed triggers from update 15 to 65. The matched permutation preserves the distribution within class.}{statetrigger}
\tab{Motion quality under state-triggered GAF scoring versus matched terminal selection. Differences are state-triggered minus terminal. Intervals resample the eight fresh episodes 2,000 times, conditional on the frozen models and draws. Stop discrepancy is in percentage points.}{statequality}{\begin{tabular}{lrrl}\toprule Metric & In-loop & Terminal & Difference [95\% interval]\\\midrule
Speed $\W$, m/s &1.1553&1.1492&$+0.0061$ [$-0.0290$, $+0.0347$]\\
Accel. $\W$, m/s$^2$ &0.2178&0.2174&$+0.0004$ [$-0.0023$, $+0.0030$]\\
Travel $\W$, m &144.98&144.01&$+0.97$ [$-1.96$, $+3.29$]\\
Stop discrepancy, pp &3.59&3.52&$+0.07$ [$-0.14$, $+0.24$]\\\bottomrule\end{tabular}}
The state rule gave a slightly larger speed discrepancy than matched terminal scoring. Permuted timing gave a slightly smaller discrepancy than its terminal comparator, $-0.0054$ m/s. The difference of these paired effects was $+0.0115$ m/s, interval $[-0.0311,0.0385]$. The experiment did not establish an advantage from adapting the intervention to motion stability. It was not an equivalence test and does not exclude small effects or effects of different feedback functions.

All 2,160 traces under each main policy passed the directed graph audit. Four traces in each main arm exhausted their finite path; all were retained in the metrics. Mean coordinate separation across independently returned draws was 58.65 metres for state triggering and 56.82 metres for terminal scoring. Thus the state test changed the allocation of continued denoising without producing a consistent improvement in the measured class-conditional motion distributions. On the present evidence, shared-road class switching is the clearest demonstration of learned motion control; image preference is an experimentally tested, optional intervention.
